%% ============================================================
%%  APS ARTIFACT TEMPLATE
%%  Example skeleton for a stylized "artifact edition" manuscript,
%%  using aps-artifact-theme.sty. Copy this file into the target
%%  volume (tslca/, ftqc/, etc.), rename it, fill in the macros
%%  below, add frame art with \SetArtifactFrames if you have it,
%%  and start writing chapters with \artifactchapter / \artifactendsection.
%%
%%  This file has NO Fuxyez-specific content — it's the generic base.
%%  Compile with XeLaTeX (fontspec requires it).
%% ============================================================

\documentclass[12pt, a4paper]{book}

%% Path is relative to wherever you copy this file to. From e.g.
%% ftqc/aps_ftqc_skrypt.tex or tslca/tslca_skrypt.tex, that's "../".
\usepackage{../aps-artifact-theme}

%% --- Optional: override theme fonts/colors here, AFTER the \usepackage line ---
% \setmainfont{Your Serif Font}
% \setmonofont{Your Mono Font}
% \definecolor{artifactAccentA}{HTML}{XXXXXX}   % re-theme the primary accent
% \definecolor{artifactAccentB}{HTML}{XXXXXX}   % re-theme the secondary accent

%% --- Optional: frame art for odd/even pages (omit for a plain background) ---
% \SetArtifactFrames{path/to/odd_frame.png}{path/to/even_frame.png}

%% --- Header / footer labels (all optional, blank if unset) ---
\SetArtifactHeaderLeft{[ EXAMPLE LABEL ]}
\SetArtifactHeaderCenter{ARTIFACT TEMPLATE}
\SetArtifactHeaderRight{[ EXAMPLE.ORG ]}
\SetArtifactFooterLeft{Aurphyx LLC -- Ross A. Edwards}
\SetArtifactFooterRight{v0.0.1}

%% --- Title page content (all optional, blank if unset) ---
\artifacttitle{EXAMPLE ARTIFACT TITLE}
\artifactsubtitle{A one-line subtitle for the work}
\SetArtifactEyebrow{~~~ optional eyebrow line above the title ~~~}
\SetArtifactTagline{OPTIONAL TAGLINE LINE\\ACROSS ONE OR MORE LINES}
\SetArtifactByline{SCRIBED BY [ AUTHOR ]\\AURPHYX LLC}
\SetArtifactEdition{v0.0.1 -- Example Edition -- example.org}
\SetArtifactClosingNote{This document is an example. Replace this note.}

%% --- Chapter/section banner numbering ---
\SetArtifactProtocolLabel{EXAMPLE PROTOCOL}
\SetArtifactTotalSections{2}

\begin{document}

\artifacttitlepage
\artifacttoc

%% Optional CLI-mockup box on the title page — call this INSIDE
%% \artifacttitlepage's flow is not supported (titlepage macro is
%% self-contained); instead, add one as its own chapter opener or
%% inline within a chapter body via \ArtifactCLIBox{...} wherever
%% you want the terminal-mockup look, e.g.:
%%
%%   \ArtifactCLIBox{
%%   \$\$ example install thing --flag=value \\
%%   > Doing the thing... \\
%%   > Done.
%%   }

\artifactchapter{First Example Chapter}

This is body text. Use standard LaTeX sectioning
(\texttt{\textbackslash section}, \texttt{\textbackslash subsection}) inside
an \texttt{\textbackslash artifactchapter} block — the theme already styled
those commands, no extra markup needed here.

\begin{artifactnote}{EXAMPLE NOTE}
This is a callout box, for asides or flagged content.
\end{artifactnote}

\begin{artifactglyphbox}
This is the plainer terminal/banner box, useful for quoted fragments
or short emphasized text blocks.
\end{artifactglyphbox}

\begin{lstlisting}
// Code listings inherit the theme's dark background and line numbers.
// Call \lstdefinelanguage{yourlang}{...} before \begin{document} and
// \lstset{language=yourlang} to get keyword highlighting.
fn example() -> Result {
    return ok;
}
\end{lstlisting}

\artifactendsection

\artifactchapter{Second Example Chapter}

More body content here.

\artifactendsection

\end{document}
